% GENERATED FILE. Edit canonical lessons, then run npm run generate:books.
% canonical-source-hash: fnv1a64:5dcdc1cc2dd00429
% canonical-lessons: IT-C444-quartiere, IT-C444-periferia, IT-C444-rintracciare, IT-C444-numero-civico, IT-C444-citofono, IT-R444-ripasso-437-444

\chapter{L'indirizzo}
\label{ch:it-l-indirizzo}
\hlchaptermodality{444}

\begin{chapteropening}
\textbf{By the end of this chapter:} \emph{I can find an address: the neighbourhood, the outskirts, tracking someone down, the house number and the intercom.}

Complete the last lesson of chapter 444: I can find an address: the neighbourhood, the outskirts, tracking someone down, the house number and the intercom.
\end{chapteropening}

\section[il quartiere]{il quartiere — the neighbourhood, the district}

\label{lesson:IT-C444-quartiere}



\begin{quote}
Before the new one: say the Italian for western, then the Italian for the landmark.
\end{quote}

\subsection*{You'll want to know: il quartiere}
\textbf{il quartiere} — \textquotedblleft{}the neighbourhood, the district\textquotedblright{}.

\begin{quote}
Abitiamo in un quartiere tranquillo.
\end{quote}

\begin{quote}
\emph{We live in a quiet neighbourhood.}
\end{quote}

\emph{Il quartier generale} — headquarters.

\begin{cousinweb}
\emph{quartiere} was borrowed from French \emph{quartier} (EtymDB), from Latin \emph{quartārius} 'a fourth part': first a quarter of a town.
\end{cousinweb}

\begin{grammarlens}[title={What you've built}]
One more word for telling, explaining and describing at B1.
\end{grammarlens}

\subsection*{Your turn}
\begin{itemize}
\raggedright
  \item \emph{Say it:} \emph{il quartiere}
  \item \emph{Say it:} \emph{il quartiere}, once more
  \item \emph{Say it:} a sentence of your own with \emph{il quartiere}
  \item \emph{Recall:} say the Italian for western, then the Italian for the landmark, then say \emph{il quartiere} again
\end{itemize}

\begin{tcolorbox}[breakable,colback=teal!4,colframe=teal!35!black,title={Before you move on}]
What does il quartiere mean? (\textquotedblleft{}The neighbourhood, the district\textquotedblright{}.) Say it once more.
\end{tcolorbox}
\section[la periferia]{la periferia — the outskirts, the suburbs}

\label{lesson:IT-C444-periferia}



\begin{quote}
Before the new one: say the Italian for the landmark, then the Italian for the neighbourhood.
\end{quote}

\subsection*{You'll want to know: la periferia}
\textbf{la periferia} — \textquotedblleft{}the outskirts, the suburbs\textquotedblright{}.

\begin{quote}
La casa è in periferia, lontano dal centro.
\end{quote}

\begin{quote}
\emph{The house is on the outskirts, far from the centre.}
\end{quote}

Often less smart than the English suburbs.

\begin{cousinweb}
\emph{periferia} is a learned borrowing from Late Latin \emph{peripherīa}, from Greek \emph{periphéreia} 'circumference', built on \emph{pherein} 'to carry'.
\end{cousinweb}

\begin{grammarlens}[title={What you've built}]
One more word for telling, explaining and describing at B1.
\end{grammarlens}

\subsection*{Your turn}
\begin{itemize}
\raggedright
  \item \emph{Say it:} \emph{la periferia}
  \item \emph{Say it:} \emph{la periferia}, once more
  \item \emph{Say it:} a sentence of your own with \emph{la periferia}
  \item \emph{Recall:} say the Italian for the landmark, then the Italian for the neighbourhood, then say \emph{la periferia} again
\end{itemize}

\begin{tcolorbox}[breakable,colback=teal!4,colframe=teal!35!black,title={Before you move on}]
What does la periferia mean? (\textquotedblleft{}The outskirts, the suburbs\textquotedblright{}.) Say it once more.
\end{tcolorbox}
\section[rintracciare]{rintracciare — to track down, to trace}

\label{lesson:IT-C444-rintracciare}



\begin{quote}
Before the new one: say the Italian for the neighbourhood, then the Italian for the outskirts.
\end{quote}

\subsection*{You'll want to know: rintracciare}
\textbf{rintracciare} — \textquotedblleft{}to track down, to trace\textquotedblright{}.

\begin{quote}
Non riesco a rintracciare il suo indirizzo.
\end{quote}

\begin{quote}
\emph{I cannot track down his address.}
\end{quote}

\emph{Essere rintracciabile} — to be reachable.

\begin{cousinweb}
\emph{rintracciare} is built on \emph{traccia} 'trace, track', the noun of \emph{tracciare} 'to trace', which is usually traced to Vulgar Latin \emph{tractiāre}, built on \emph{trahere} 'to drag' (Lewis \& Short).
\end{cousinweb}

\begin{grammarlens}[title={What you've built}]
One more word for telling, explaining and describing at B1.
\end{grammarlens}

\subsection*{Your turn}
\begin{itemize}
\raggedright
  \item \emph{Say it:} \emph{rintracciare}
  \item \emph{Say it:} \emph{rintracciare}, once more
  \item \emph{Say it:} a sentence of your own with \emph{rintracciare}
  \item \emph{Recall:} say the Italian for the neighbourhood, then the Italian for the outskirts, then say \emph{rintracciare} again
\end{itemize}

\begin{tcolorbox}[breakable,colback=teal!4,colframe=teal!35!black,title={Before you move on}]
What does rintracciare mean? (\textquotedblleft{}To track down, to trace\textquotedblright{}.) Say it once more.
\end{tcolorbox}
\section[il numero civico]{il numero civico — the house number, the street number}

\label{lesson:IT-C444-numero-civico}



\begin{quote}
Before the new one: say the Italian for the outskirts, then the Italian for to track down.
\end{quote}

\subsection*{You'll want to know: il numero civico}
\textbf{il numero civico} — \textquotedblleft{}the house number, the street number\textquotedblright{}.

\begin{quote}
Qual è il numero civico?
\end{quote}

\begin{quote}
\emph{What is the house number?}
\end{quote}

Written after the street name: \emph{via del Mare, 12}.

\begin{cousinweb}
\emph{numero civico} is built from \emph{numero} 'number', inherited from Latin \emph{numerus}, and \emph{civico}, a learned borrowing from Latin \emph{cīvicus} 'of citizens', from \emph{cīvis} 'citizen' (EtymDB).
\end{cousinweb}

\begin{grammarlens}[title={What you've built}]
One more word for telling, explaining and describing at B1.
\end{grammarlens}

\subsection*{Your turn}
\begin{itemize}
\raggedright
  \item \emph{Say it:} \emph{il numero civico}
  \item \emph{Say it:} \emph{il numero civico}, once more
  \item \emph{Say it:} a sentence of your own with \emph{il numero civico}
  \item \emph{Recall:} say the Italian for the outskirts, then the Italian for to track down, then say \emph{il numero civico} again
\end{itemize}

\begin{tcolorbox}[breakable,colback=teal!4,colframe=teal!35!black,title={Before you move on}]
What does il numero civico mean? (\textquotedblleft{}The house number, the street number\textquotedblright{}.) Say it once more.
\end{tcolorbox}
\section[il citofono]{il citofono — the intercom, the entry phone}

\label{lesson:IT-C444-citofono}



\begin{quote}
Before the new one: say the Italian for to track down, then the Italian for the house number.
\end{quote}

\subsection*{You'll want to know: il citofono}
\textbf{il citofono} — \textquotedblleft{}the intercom, the entry phone\textquotedblright{}.

\begin{quote}
Suona al citofono, ti apro.
\end{quote}

\begin{quote}
\emph{Ring the intercom, I will let you in.}
\end{quote}

\emph{Citofonare} — to buzz someone.

\begin{cousinweb}
\emph{citofono} is a learned formation, usually explained as Latin \emph{citō} 'quickly' (Lewis \& Short) and Greek \emph{phōn\'{\={e}}} 'voice'.
\end{cousinweb}

\begin{grammarlens}[title={What you've built}]
One more word for telling, explaining and describing at B1.
\end{grammarlens}

\subsection*{Your turn}
\begin{itemize}
\raggedright
  \item \emph{Say it:} \emph{il citofono}
  \item \emph{Say it:} \emph{il citofono}, once more
  \item \emph{Say it:} a sentence of your own with \emph{il citofono}
  \item \emph{Recall:} say the Italian for to track down, then the Italian for the house number, then say \emph{il citofono} again
\end{itemize}

\begin{tcolorbox}[breakable,colback=teal!4,colframe=teal!35!black,title={Before you move on}]
What does il citofono mean? (\textquotedblleft{}The intercom, the entry phone\textquotedblright{}.) Say it once more.
\end{tcolorbox}
\section[(dialogue)]{Ripasso dei capitoli 437–444}

\label{lesson:IT-R444-ripasso-437-444}



\begin{quote}
Say the Italian for the racket and the din.
\end{quote}

\subsection*{Your turn}
Say these aloud:

\begin{itemize}
\raggedright
  \item the racket, the din, the cockroach, the mould, to creak
  \item the complaint, to complain, the grumble, to put right, to compensate
  \item to find one's way, to get lost, the map, the street plan, the compass
  \item the satnav, the direction, to turn, to carry on, to take
  \item the fork, the alley, the side street, the avenue, the one-way street
  \item the climb, the descent, to run along, to come out, to go round
  \item northern, southern, eastern, western, the landmark
  \item the neighbourhood, the outskirts, to track down, the house number, the intercom
\end{itemize}

\begin{tcolorbox}[breakable,colback=teal!4,colframe=teal!35!black,title={Before you move on}]
The racket? (\textbf{il baccano}.) The din? (\textbf{il chiasso}.) The cockroach? (\textbf{lo scarafaggio}.) The mould? (\textbf{la muffa}.) To creak? (\textbf{cigolare}.) The complaint? (\textbf{il reclamo}.) To complain? (\textbf{reclamare}.) The grumble? (\textbf{la lamentela}.) To put right? (\textbf{rimediare}.) To compensate? (\textbf{risarcire}.) To find one's way? (\textbf{orientarsi}.) To get lost? (\textbf{smarrirsi}.) The map? (\textbf{la cartina}.) The street plan? (\textbf{la piantina}.) The compass? (\textbf{la bussola}.) The satnav? (\textbf{il navigatore}.) The direction? (\textbf{l'indicazione}.) To turn? (\textbf{svoltare}.) To carry on? (\textbf{proseguire}.) To take? (\textbf{imboccare}.) The fork? (\textbf{il bivio}.) The alley? (\textbf{il vicolo}.) The side street? (\textbf{la traversa}.) The avenue? (\textbf{il viale}.) The one-way street? (\textbf{il senso unico}.) The climb? (\textbf{la salita}.) The descent? (\textbf{la discesa}.) To run along? (\textbf{costeggiare}.) To come out? (\textbf{sbucare}.) To go round? (\textbf{aggirare}.) Northern? (\textbf{settentrionale}.) Southern? (\textbf{meridionale}.) Eastern? (\textbf{orientale}.) Western? (\textbf{occidentale}.) The landmark? (\textbf{il punto di riferimento}.) The neighbourhood? (\textbf{il quartiere}.) The outskirts? (\textbf{la periferia}.) To track down? (\textbf{rintracciare}.) The house number? (\textbf{il numero civico}.) The intercom? (\textbf{il citofono}.)
\end{tcolorbox}
